\section*{迷向元素与奇异元素：}

设 \(\Phi\) 是左 A-模 E 上的一个 \(\varepsilon\)-埃尔米特型。称元素 \(x \in E\) 为迷向的，如果 \(\Phi(x, x) = 0\) ；称 E 的子模 M 为迷向的，如果 \(\mathbf{M} \cap \mathbf{M}^{\circ} \neq \{0\}\) （换言之，如果 \(\Phi\) 在 M 上的限制是退化的）；称 M 为全迷向的，如果 \(\mathbf{M} \subset \mathbf{M}^{\circ}\) （换言之，如果 \(\Phi\) 在 M 上的限制为零）。对 E 的每个子模 M，\(\mathbf{M} \cap \mathbf{M}^{\circ}\) 都是全迷向的。若 E 是有限维向量空间且 \(\Phi\) 非退化，则对 E 的子空间 M，下列三个条件等价： \(1^{\circ}\) M 非迷向； \(2^{\circ}\) M \(^{\circ}\) 非迷向； \(3^{\circ}\) E 是 M 与 M \(^{\circ}\) 的直和。设环 A 是交换环，Q 是 A-模 E 上的一个二次型，\(\Phi\) 是与 Q 相伴的双线性型。称元素 \(x \in E\) 为奇异的，如果 \(\mathrm{Q}(x) = 0\) ；称 E 的子模 M 为奇异的（相应地，全奇异的），如果 \(\mathbf{M} \cap \mathbf{M}^{\circ}\) 含有一个 \(\neq 0\) 的奇异元素（相应地，如果 Q 在 M 上的限制为零）。每个奇异元素（相应地，奇异子模、全奇异子模）都是迷向的（相应地，迷向的、全迷向的）；当 A 是特征 \(\neq 2\) 的域时，其逆亦真。
